\subsection{Fixed-dimensional support recovery}
\label{sec:higher-dimensional-support-counts}

This subsection proves the fixed-dimensional support identity used in both
parts of Theorem~\ref{thm:main-support-law}. Exposed radial points are
separated by much more than the \(1/\rho_n\) width of a likelihood patch, so
their local optimizations decouple: a superthreshold patch contains one fitted
atom and a subthreshold patch contains none. Here
\(\rho_n=\sqrt{2\log n}\) for \(d=2\) and
\(\rho_n=\sqrt{2b_{n,d}}\) for fixed \(d\ge3\), while \(\tau\) denotes the
central score location from Lemma~\ref{lem:support-fixed-core-collapse}. Put
\(\Delta_n^{\rm anc}=\gamma_n\tau-\bar Z_n\). The exact anchor identity is
\begin{equation}
\label{eq:support-high-d-anchor-translation}
 D_{n,\tau}(a)
 =\frac1n\sum_iq_{n,\tau,a}(i)
 =e^{-a\cdot\Delta_n^{\rm anc}}D_{n,*}(a),
 \qquad
 D_{n,*}(a)=\frac1n\sum_i
 e^{a\cdot(Z_i-\bar Z_n)-\gamma_n|a|^2/2}.
\end{equation}
The central location is tied to the sample mean by an exact barycenter
identity. For a score-coordinate NPMLE \(G\), let
\[
 \Pi_i(dt)=
 \frac{\exp\{t\cdot Z_i-\gamma_n|t|^2/2\}\,G(dt)}{m_G(Z_i)}
\]
be the posterior mixing distribution for observation \(i\). KKT equality
\(G\)-almost everywhere gives
\[
 \frac1n\sum_{i=1}^n\Pi_i(dt)=G(dt).
\]
On the other hand, differentiating the likelihood when every score location is
translated by the same vector gives
\[
 0=\frac1n\sum_{i=1}^n
 \left\{Z_i-\gamma_n\int t\,\Pi_i(dt)\right\}.
\]
Consequently,
\begin{equation}
\label{eq:support-npmle-barycenter}
    \bar Z_n=\gamma_n\int t\,G(dt).
\end{equation}
For the decomposition in Lemma~\ref{lem:support-fixed-core-collapse}, this
implies
\[
 \Delta_n^{\rm anc}
 =\gamma_n\int(\tau-t)\,G_n^{\rm out}(dt).
\]
Support reduction places \(G\) in
\(\operatorname{conv}\{Z_1/\gamma_n,\ldots,Z_n/\gamma_n\}\). Since
\(\max_i|Z_i|=O_{\Pbb}(\rho_n)\) in fixed dimension, the rate in
Lemma~\ref{lem:support-fixed-core-collapse} therefore gives
\begin{equation}
\label{eq:support-high-d-anchor-error}
    \rho_n|\Delta_n^{\rm anc}|
    =O_{\Pbb}(W_n\rho_n^2)
    =o_{\Pbb}(\kappa_n)
\end{equation}
in both fixed-dimensional critical windows. For the last equality, use
\[
 \rho_n^2=O(\log n),\qquad
 W_n=O_{\Pbb}\left(\frac{(\log n)^{C_d}}{n\kappa_n}\right),
 \qquad \kappa_n\gtrsim\frac1{\log n}.
\]

We use one symbol for the dimension-specific radial excess:
\begin{equation}
\label{eq:fixed-support-radial-excess}
 x_{n,i}=
 \begin{cases}
   |Z_i|^2/2-\log n,&d=2,\\[2mm]
   |Z_i-\bar Z_n|^2/2-b_{n,d},&d\ge3.
 \end{cases}
\end{equation}
In fixed dimension, replacing the centered radius in the second line by
\(|Z_i|\) changes all excesses by \(o_{\Pbb}(1)\), uniformly in \(i\).

We now fix the local coordinate maps used below. In \(d=2\), write
\(U_j=Z_j/|Z_j|\), let \(\mathsf R_\theta\) denote planar rotation through
angle \(\theta\), and set
\begin{equation}
\label{eq:support-d2-patch-map}
 a_{n,j}^{(2)}(r,z)
 = (\rho_n+r)\mathsf R_{z/\rho_n}U_j,
 \qquad (r,z)\in\mathbb R^2.
\end{equation}
For fixed \(d\ge3\), set
\begin{equation}
\label{eq:support-dge3-patch-map}
 a_{n,j}^{(\ge3)}(h)
 =a_{n,j}^*+h,
 \qquad
 a_{n,j}^*=\frac{Z_j-\gamma_n\tau}{\gamma_n},
 \qquad h\in\mathbb R^d.
\end{equation}
If \(\nu_j\) is a measure in the corresponding local coordinates, its
pushforward under \(a_{n,j}^{(d)}\) is a measure on score displacements. Below
we write
\begin{equation}
\label{eq:support-patch-pushforward}
    \nu=\sum_j(a_{n,j}^{(d)})_\#\nu_j
\end{equation}
for the combined displacement measure.

For a finite positive measure \(\nu\) on displacement space, with
\(m=\nu(\mathbb R^d)\le n\), consider the anchored mixture
\[
    G_{n,\tau,\nu}
    =\left(1-\frac mn\right)\delta_\tau
      +\frac1n(\tau+\cdot)_\#\nu .
\]
Its exact likelihood gain over \(\delta_\tau\) is
\begin{equation}
\label{eq:support-anchored-likelihood}
    \mathcal L_{n,\tau}(\nu)
    =\sum_{i=1}^n
      \log\left\{1+F_{n,i}(\nu)-\frac mn\right\},
    \qquad
    F_{n,i}(\nu)=\frac1n\int q_{n,\tau,a}(i)\,\nu(da).
\end{equation}
Thus \(m=nW\) is the natural rescaled edge mass.

\subsubsection{Local likelihood programs}

The field limits in Sections~\ref{sec:d2} and
\ref{sec:fixed-d-patches} determine the likelihood program associated with one
radial extreme. Optimizing this program shows that an active patch contains one
atom rather than a split mixing measure.

Suppose first that \(d\ge3\), along a sequence on which \(s_n\to s\). For an
observation of centered radial excess \(x\), completing the square at the
displacement \(a=a_{n,i}^{(\ge3)}(h)\) gives
\[
    \frac1nq_{n,\tau,a}(i)\longrightarrow
    A_x(h):=e^{x-s-|h|^2/2}.
\]
If a finite positive measure \(\nu\) of rescaled mass is placed in this
patch, its limiting gain is
\begin{equation}
\label{eq:support-dge3-local-program}
    \mathcal J_x^{(\ge3)}(\nu)
    =\log\left\{1+\int A_x(h)\,\nu(dh)\right\}
      -\nu(\mathbb R^d).
\end{equation}
For fixed total mass, the integral is uniquely maximized by putting all mass
at \(h=0\). The remaining scalar problem is
\[
    \max_{y\ge0}\{\log(1+e^{x-s}y)-y\},
\]
whose optimizer is \(y_x^*=(1-e^{-(x-s)})_+\). Hence the patch is active
exactly when \(x>s\), and every nonzero limiting optimizer is the one-atom
measure \(y_x^*\delta_0\).

In dimension two, use radial score offset \(r\) and blown-up angular
displacement \(z\), and write
\[
    B_\lambda(r)=e^{-\lambda}\Phi(-r),
    \qquad
    C_{\lambda,x}(r,z)=e^{-\lambda}e^{x-r^2/2-z^2/2}.
\]
The corresponding local program is
\begin{equation}
\label{eq:support-d2-local-program}
    \mathcal J_x^{(2)}(\nu)
    =\log\left\{1+\int C_{\lambda,x}\,d\nu\right\}
      -\int\{1-B_\lambda(r)\}\,d\nu(r,z).
\end{equation}
Put
\[
    R_x(r,z)=\frac{C_{\lambda,x}(r,z)}{1-B_\lambda(r)}.
\]
If \(u=\int C_{\lambda,x}\,d\nu\), then
\[
    \int(1-B_\lambda)\,d\nu\ge
    \frac{u}{\sup_{r,z}R_x(r,z)},
\]
with equality only when \(\nu\) is supported on the maximizers of \(R_x\).
The angular maximizer is uniquely \(z=0\). The radial maximizer is uniquely
\(r=-a_\lambda\), where
\[
    e^\lambda=\Phi(a_\lambda)+\frac{\phi(a_\lambda)}{a_\lambda},
\]
and the calculation in Section~\ref{sec:d-two-sharp-window} gives
\(\sup R_x=e^{x-x_\lambda}\). Thus the patch is active exactly when
\(x>x_\lambda\), and every nonzero limiting optimizer is a single atom at
\((-a_\lambda,0)\). Its rescaled mass is
\[
    q_x^*=
    \frac{1-e^{-(x-x_\lambda)}}
         {1-e^{-\lambda}\Phi(a_\lambda)}.
\]

For fixed \(d\ge3\), the contact argument requires differentiated control near
a random radial extreme.  The next lemma supplies that control without
conditioning on a sample-centered exceedance event.  It first fixes the raw
extreme, so the punctured sample remains independent, and then sums over the
possible distinguished observations by the Palm formula.
\begin{lemma}
\label{lem:support-dge3-palm-c2}
Fix \(d\ge3\) and suppose \(s_n=O(1)\).  Put
\[
 \widetilde x_{n,i}=\frac{|Z_i|^2}{2}-b_{n,d},
 \qquad
 \mathcal E_n(C)=\{j:|\widetilde x_{n,j}|\le C\}.
\]
For fixed \(A,C,R<\infty\), define, for \(j\in\mathcal E_n(C)\),
\[
 \mathcal R_{n,j,A}(h)
 =\frac1n\sum_{i\ne j}
   \exp\left\{
      \left(\frac{Z_j}{\gamma_n}+h\right)\!\cdot Z_i
      -\frac{\gamma_n}{2}
       \left|\frac{Z_j}{\gamma_n}+h\right|^2
   \right\}
   \mathbf 1_{\{\widetilde x_{n,i}\le-A\}},
 \qquad h\in B_R.
\]
Then
\begin{equation}
\label{eq:support-dge3-palm-c2}
 \max_{j\in\mathcal E_n(C)}
 \|\mathcal R_{n,j,A}\|_{C^2(B_R)}\xrightarrow{\Pbb}0.
\end{equation}
Moreover, if \(K_{n,j,i}(h)\) denotes the exponential summand in the
definition of \(\mathcal R_{n,j,A}\), then
\begin{equation}
\label{eq:support-dge3-palm-absolute}
 \max_{j\in\mathcal E_n(C)}\max_{|\beta|\le2}\sup_{h\in B_R}
 \frac1n\sum_{i\ne j}
 |\partial_h^\beta K_{n,j,i}(h)|
 \mathbf1_{\{\widetilde x_{n,i}\le-A\}}
 \xrightarrow{\Pbb}0.
\end{equation}
The same conclusions hold with the candidate condition
\(|x_{n,j}|\le C\) and lower-cloud cutoff \(x_{n,i}\le-A\), where \(x_{n,i}\)
is the sample-centered excess in \eqref{eq:fixed-support-radial-excess}.
\end{lemma}

\begin{proof}
Let \(R_{n,A}^2/2=b_{n,d}-A\), choose \(p>d\), and fix a deterministic
\(z\) with \(\bigl||z|^2/2-b_{n,d}\bigr|\le C\).  For
\(t_z(h)=z/\gamma_n+h\), set
\[
 V_{n,z,i}(h)=\frac1n
 e^{t_z(h)\cdot Z_i-\gamma_n|t_z(h)|^2/2}
 \mathbf 1_{\{|Z_i|\le R_{n,A}\}}.
\]
We first establish, uniformly over these \(z\), \(h\in B_R\),
\(q\in\{2,p\}\), and \(|\beta|\le3\),
\begin{equation}
\label{eq:support-palm-derivative-moments}
 \sum_{i=2}^n
 \E|\partial_h^\beta V_{n,z,i}(h)|^q
 \le C_{A,C,R,p,d}\rho_n^{-1}e^{-\lambda_n}.
\end{equation}
Indeed, differentiation inserts a polynomial of degree at most \(|\beta|\)
in \(Z_i-\gamma_nt_z(h)\).  Exponential tilting therefore gives
\begin{align*}
 &\sum_{i=2}^n
 \E|\partial_h^\beta V_{n,z,i}(h)|^q\\
 &\quad\le
 C n^{1-q}e^{-q\kappa_n|t_z(h)|^2+(q^2-q)|t_z(h)|^2/2}
 \E\left[
  \{1+|G+(q-\gamma_n)t_z(h)|\}^{q|\beta|}
  \mathbf1_{\{|G+qt_z(h)|\le R_{n,A}\}}
 \right].
\end{align*}
Aligning \(t_z(h)=re_1\), conditioning on \(G_\perp\), and applying Mills'
bound to \(G_1\) show that the last expectation is at most
\[
 C_{A,C,R,p,d}\rho_n^{-1}
   \exp\{-(qr-R_{n,A})^2/2\}.
\]
Here the polynomial causes only a constant loss: on the truncation event its
longitudinal argument is
\(R_{n,A}-\gamma_nr+O(1/\rho_n)\), which is bounded for
\(|r-|z|/\gamma_n|\le R\), while the transverse Gaussian moments are fixed.
Equivalently, this weighted Mills estimate follows by integrating
\((1+|u|)^m e^{-u^2/2}\) beyond the same one-dimensional boundary.

The exponent outside \(\rho_n^{-1}\) equals
\begin{align*}
 &(1-q)\log n-q\kappa_nr^2
   +\frac{q^2-q}{2}r^2-\frac12(qr-R_{n,A})^2\\
 &\qquad=(q-1)\{b_{n,d}-\log n-A\}
   -q\kappa_nr^2-\frac q2(r-R_{n,A})^2.
\end{align*}
Since \(b_{n,d}-\log n=\lambda_n-s_n\),
\(\kappa_nr^2=\lambda_n+o(1)\), and \(s_n=O(1)\), it is at most
\(-\lambda_n+O_{A,C,R,p,d}(1)\).  This proves
\eqref{eq:support-palm-derivative-moments}.  The same tilting calculation with
\(q=1\) gives
\begin{equation}
\label{eq:support-palm-derivative-means}
 \sup_{z,h,|\beta|\le3}
 \left|\sum_{i=2}^n\E\partial_h^\beta V_{n,z,i}(h)\right|
 \le C_{A,C,R,d}e^{-\lambda_n}=o(1).
\end{equation}
The same \(q=1\) calculation applies to the absolute values of the
derivatives.  Using the derivatives of order three and the fundamental theorem
of calculus on a fixed mesh of \(B_R\) gives
\[
 \sup_z\E\left[
  \max_{|\beta|\le2}\sup_{h\in B_R}
  \frac1n\sum_{i=2}^n
  |\partial_h^\beta K_{n,z,i}(h)|
  \mathbf1_{\{|Z_i|\le R_{n,A}\}}
 \right]
 \le C_{A,C,R,d}e^{-\lambda_n}=o(1).
\]
The Palm union bound below and Markov's inequality therefore also prove
\eqref{eq:support-dge3-palm-absolute}.

Apply Rosenthal's inequality to the centered derivatives and integrate over
\(B_R\).  Equations \eqref{eq:support-palm-derivative-moments} and
\eqref{eq:support-palm-derivative-means} yield, uniformly in \(z\),
\[
 \E\left\|
   \sum_{i=2}^nV_{n,z,i}
 \right\|_{W^{3,p}(B_R)}^p=o(1).
\]
Because \(p>d\), the embedding \(W^{3,p}(B_R)\hookrightarrow C^2(B_R)\)
and Markov's inequality imply
\[
 \sup_{z:\,\lvert |z|^2/2-b_{n,d}\rvert\le C}
 \Pbb\left\{
  \left\|\sum_{i=2}^nV_{n,z,i}\right\|_{C^2(B_R)}>\delta
 \right\}=o(1)
\]
for every \(\delta>0\).  Independence now gives the Palm union bound
\begin{align*}
 &\Pbb\left\{
   \max_{j\in\mathcal E_n(C)}
   \|\mathcal R_{n,j,A}\|_{C^2(B_R)}>\delta
 \right\}\\
 &\quad\le n\int_{\{\lvert |z|^2/2-b_{n,d}\rvert\le C\}}
  \Pbb\left\{
   \left\|\sum_{i=2}^nV_{n,z,i}\right\|_{C^2(B_R)}>\delta
  \right\}\,P(dz)=o(1),
\end{align*}
because the Gamma tail estimate gives
\(nP\{\lvert |Z|^2/2-b_{n,d}\rvert\le C\}=O_C(1)\).

Finally,
\[
 \max_i\left|
  \frac{|Z_i-\bar Z_n|^2-|Z_i|^2}{2}
 \right|=o_{\Pbb}(1).
\]
Fix \(\delta>0\).  Apart from an event of probability tending to zero,
\(|x_{n,j}|\le C\) implies \(|\widetilde x_{n,j}|\le C+\delta\), while
\(x_{n,i}\le-A\) implies
\(\widetilde x_{n,i}\le-(A-\delta)\).  Applying the raw result with
candidate bound \(C+\delta\) and lower cutoff \(A-\delta\) proves the
sample-centered assertion directly.
\end{proof}

\subsubsection{Uniform patch stability}

We need the local approximations uniformly over the random edge patches and
all candidate mixing measures. As in
Sections~\ref{sec:sharp-high-dimensional-estimates} and
\ref{sec:d-two-sharp-window}, the lower radial cutoff is removed after the
finite-sample limit.

The candidate patches and the observational cutoff play different roles.  The
candidate set is fixed by the activation threshold and remains unchanged as
the cutoff is removed.  The cutoff is used only to divide the data into a
finite exposed cloud and a lower cloud.  The following estimate controls the
lower cloud uniformly over every mixing measure on the candidate patches.
\begin{lemma}
\label{lem:support-uniform-lower-cloud}
Fix \(C,R,M,J<\infty\).  Let \(\mathcal I_n\) be any set of at most \(J\)
candidate observations with \(|x_{n,j}|\le C\), and let \(\nu_j\) be supported
in \(B_R\) for \(d\ge3\), or in \([-R,R]^2\) for \(d=2\), with
\(m=\sum_{j\in\mathcal I_n}\nu_j(\mathbb R^d)\le M\).  Form \(\nu\) by
\eqref{eq:support-patch-pushforward} and write
\(F_{n,i}=F_{n,i}(\nu)\).  Suppose \(s_n=O(1)\) for fixed \(d\ge3\), or
\(\lambda_n\to\lambda\in(0,\infty)\) for \(d=2\).

There is \(A_0=A_0(C,R,M)>C+1\) such that, for each fixed
\(A\ge A_0\), on the event \(\max_i x_{n,i}\le C\) and the usual
edge-separation event,
\begin{equation}
\label{eq:support-lower-cloud-max-square}
 \max_{i:x_{n,i}\le-A}F_{n,i}\le C_{C,R}Me^{-A},
 \qquad
 \sum_{i:x_{n,i}\le-A}F_{n,i}^2
 \le C_{C,R}M^2e^{-A}+o_{\Pbb,A}(1),
\end{equation}
apart from an event of probability tending to zero.  For every candidate
patch, one also has the derivative-envelope bound
\begin{equation}
\label{eq:support-lower-cloud-derivative-envelope}
 \max_{|\beta|\le2}\sup_u
 \frac1n\sum_{i:x_{n,i}\le-A}
 \left|\partial_u^\beta
 q_{n,\tau,a_{n,j}(u)}(i)\right|
 =O_{\Pbb,A_0,C,R}(1),
\end{equation}
uniformly for \(A\ge A_0\).
Uniformly over the same measures,
\begin{equation}
\label{eq:support-lower-cloud-linear}
 \sum_{i:x_{n,i}\le-A}F_{n,i}
 =
 \begin{cases}
   o_{\Pbb,A}(1),&d\ge3,\\[1mm]
   \displaystyle
   \sum_{j\in\mathcal I_n}\int B_{\lambda_n}(r)\,\nu_j(dr,dz)
      +O_{\Pbb,C,R,J}(Me^{-A})+o_{\Pbb,A}(1),&d=2,
 \end{cases}
\end{equation}
where
\[
 B_{\lambda_n}(r)=e^{-\lambda_n}\Phi(-r).
\]
The observations with \(-A<x_{n,i}\le C\) that do not belong to
\(\mathcal I_n\) contribute \(o_{\Pbb,A}(1)\), uniformly over the candidate
patches and measures.
\end{lemma}

\begin{proof}
Completing the square gives, for every score displacement \(a\),
\begin{equation}
\label{eq:support-lower-cloud-square-completion}
 \frac1nq_{n,\tau,a}(i)
 =\exp\left\{
   \frac{|Z_i-\gamma_n\tau|^2}{2\gamma_n}-\log n
   -\frac{\gamma_n}{2}
      \left|a-\frac{Z_i-\gamma_n\tau}{\gamma_n}\right|^2
 \right\}.
\end{equation}
For \(d\ge3\), the anchor estimate
\eqref{eq:support-high-d-anchor-error} and
\(b_{n,d}-\log n=\lambda_n-s_n\) make the first exponent at most
\(-A-s_n+o_{\Pbb}(1)\) when \(x_{n,i}\le-A\).  For \(d=2\), the radial
identity in Lemma~\ref{lem:support-d2-c2-absorption} bounds the exponent by
\(-A+O_{C,R}(1)\). Integrating
\eqref{eq:support-lower-cloud-square-completion} over a measure of mass at most
\(M\) proves the first inequality in
\eqref{eq:support-lower-cloud-max-square}.

For fixed \(d\ge3\), Lemma~\ref{lem:support-dge3-palm-c2}, followed by the
translation from \(Z_j/\gamma_n\) to
\((Z_j-\gamma_n\tau)/\gamma_n\), gives
\[
 \max_{j\in\mathcal I_n}\sup_{u}
 \frac1n\sum_{i:x_{n,i}\le-A}
 q_{n,\tau,a_{n,j}^{(\ge3)}(u)}(i)=o_{\Pbb,A}(1).
\]
The translation is legitimate in \(C^2\) because
\(|\tau|=o_{\Pbb}(1)\) and the Palm lemma controls a slightly larger compact
set.  The translated Palm field is multiplied by
\(e^{-\gamma_na\cdot\tau}\); uniformly on these patches its logarithm is
\(O_{\Pbb}(\rho_n|\tau|)=o_{\Pbb}(1)\) by
Lemma~\ref{lem:support-fixed-core-collapse}.  Positivity and Fubini now give the first line of
\eqref{eq:support-lower-cloud-linear}.
Equation \eqref{eq:support-dge3-palm-absolute} gives
\eqref{eq:support-lower-cloud-derivative-envelope} in this case.  The bound is
uniform for \(A\ge A_0\), because the absolute derivative sum is nonincreasing
as the lower cutoff \(-A\) is moved inward.

For \(d=2\), use the fixed-cutoff decomposition
Lemma~\ref{lem:d2-fixed-cutoff-decomposition}, the quantitative absorption
bound in the proof of Lemma~\ref{lem:d2-low-atom-absorption}, and the anchor
multiplier
\[
 e^{-a\cdot\Delta_n^{\rm anc}-a\cdot\bar Z_n}
 =e^{-\gamma_na\cdot\tau}=1+o_{\Pbb}(1)
\]
uniformly on the candidate patches. Together, these three estimates show that the lower field in a patch
with radial coordinate \(r\) is
\(B_{\lambda_n}(r)+O_{\Pbb,C,R}(e^{-A})+o_{\Pbb,A}(1)\), uniformly in its
direction.  Integrating over \(\nu\) proves the second line of
\eqref{eq:support-lower-cloud-linear}.
The radial identity bounds the absolute value of each derivative by a fixed
polynomial times the Gaussian angular kernel in
Lemma~\ref{lem:support-d2-c2-absorption}. For the deep part
\(\xi_i<-\ell_n^{\rm edge}\), repeat the boxwise Rosenthal--Morrey argument in
Lemma~\ref{lem:support-d2-c2-absorption} with a polynomial majorant for the
absolute derivatives; its mean is uniformly bounded by exponential tilting.
For the strip
\(-\ell_n^{\rm edge}\le\xi_i\le-A_0\), the slab-and-arc estimate in that
lemma, with the summable weight \((1+q^4)e^{-cq^2}\), is
\(O_{\Pbb,C,R}(e^{-A_0})\).  Their sum proves
\eqref{eq:support-lower-cloud-derivative-envelope} for \(d=2\).  Since the
absolute sum over \(\{x_{n,i}\le-A\}\) decreases with \(A\), the
\(O_{\Pbb,C,R}(e^{-A_0})\) bound holds simultaneously for all \(A\ge A_0\).

The linear bound is \(O_{\Pbb}(M)\) in either dimension.  Therefore
positivity and the first inequality in
\eqref{eq:support-lower-cloud-max-square} give
\[
 \sum_{i:x_{n,i}\le-A}F_{n,i}^2
 \le
 \left(\max_{i:x_{n,i}\le-A}F_{n,i}\right)
 \sum_{i:x_{n,i}\le-A}F_{n,i}
 \le C_{C,R}M^2e^{-A}+o_{\Pbb,A}(1),
\]
which is the second inequality in
\eqref{eq:support-lower-cloud-max-square}.

Finally, for fixed \(A,C\), the exposed cloud
\(\{-A<x_{n,i}\le C\}\) has tight cardinality.  Its distinct directions are
separated by \(\ell_n^{\rm sep}/\rho_n\), with
\(\ell_n^{\rm sep}\to\infty\), apart from an event of probability tending to
zero.  Completing the square and differentiating on fixed patches then show
that the response of every nonassociated exposed observation is
\(O(e^{-c(\ell_n^{\rm sep})^2})\).  Summing over the tight exposed cloud proves
the last assertion.
\end{proof}

Lemma~\ref{lem:support-uniform-lower-cloud} makes the lower-cloud contribution
uniform over all mixing measures on finitely many candidate patches. The next
proposition uses its linear, quadratic, and derivative bounds to reduce the
exact likelihood to a finite collection of local programs.
\begin{proposition}
\label{prop:support-finite-patch-expansion}
Fix \(C,R,M,J<\infty\), and let \(\mathcal I_n\) be any set of at most \(J\)
candidate observations with \(|x_{n,j}|\le C\).  Work on the event
\(\max_i x_{n,i}\le C\) and the edge-separation event.  For \(d\ge3\), let
\(\nu_j\) be supported in \(B_R\subset\mathbb R^d\); for \(d=2\), let it be
supported in \([-R,R]^2\).  Form the displacement measure using only
\(j\in\mathcal I_n\), assume \(\sum_{j\in\mathcal I_n}\nu_j(\mathbb R^d)\le M\),
and suppose the anchor satisfies \eqref{eq:support-high-d-anchor-error}.  Assume
that \(s_n=O(1)\) for fixed \(d\ge3\), or that
\(\lambda_n\to\lambda\in(0,\infty)\) for \(d=2\), and define
\[
\begin{aligned}
 \mathcal J_n^{(\ge3)}(\nu)
 &=\sum_{j\in\mathcal I_n}\log\left\{1+\int
 e^{x_{n,j}-s_n-|h|^2/2}\,\nu_j(dh)\right\}
 -\sum_{j\in\mathcal I_n}\nu_j(\mathbb R^d),\\
 B_{\lambda_n}(r)&=e^{-\lambda_n}\Phi(-r),\qquad
 C_{n,x}(r,z)=e^{-\lambda_n}e^{x-r^2/2-z^2/2},\\
 \mathcal J_n^{(2)}(\nu)
 &=\sum_{j\in\mathcal I_n}\log\left\{1+\int C_{n,x_{n,j}}\,d\nu_j\right\}
 -\sum_{j\in\mathcal I_n}\int\{1-B_{\lambda_n}(r)\}\,d\nu_j.
\end{aligned}
\]
There is \(A_0=A_0(C,R,M)<\infty\) such that, for every fixed
\(A\ge A_0\), uniformly over these measures,
\begin{equation}
\label{eq:support-finite-patch-error}
 \sup_\nu|\mathcal L_{n,\tau}(\nu)-\mathcal J_n^{(d)}(\nu)|
 =O_{\Pbb,C,R,J}\{(M+M^2)e^{-A}\}+o_{\Pbb,A}(1),
\end{equation}
where the random coefficient in \(O_{\Pbb,C,R,J}(1)\) is tight uniformly for
\(A\ge A_0\).  In particular, for every \(\delta>0\),
\begin{equation}
\label{eq:support-finite-patch-iterated}
 \lim_{A\to\infty}\limsup_{n\to\infty}
 \Pbb\left\{
  \sup_\nu|\mathcal L_{n,\tau}(\nu)-\mathcal J_n^{(d)}(\nu)|>\delta
 \right\}=0.
\end{equation}
The candidate space and \(\mathcal J_n^{(d)}\) do not depend on \(A\).  After
\(n\to\infty\) and then \(A\to\infty\), the limiting functional is the sum
of the independent local programs
\eqref{eq:support-dge3-local-program} or
\eqref{eq:support-d2-local-program} over the candidate patches.
\end{proposition}

\begin{proof}
Put \(m=\nu(\mathbb R^d)\), \(v_i=F_{n,i}(\nu)-m/n\), and split the data into
the candidate observations, the noncandidate exposed cloud
\(\{-A<x_{n,i}\le C\}\setminus\mathcal I_n\), and the lower cloud
\(\{x_{n,i}\le-A\}\).  For a candidate observation \(j\), completing the
square gives, uniformly on its associated patch,
\[
 \frac1nq_{n,\tau,a_{n,j}^{(\ge3)}(h)}(j)
 =e^{x_{n,j}-s_n-|h|^2/2+o(1)}
\]
in \(d\ge3\), and gives \(C_{n,x_{n,j}}(r,z)+o(1)\) in \(d=2\).
The separation event makes every cross-patch response \(o(1)\).  Hence the
candidate observations contribute the logarithmic terms in
\(\mathcal J_n^{(\ge3)}\) or \(\mathcal J_n^{(2)}\), uniformly over \(\nu\),
up to \(o_{\Pbb,A}(1)\).  The final assertion of
Lemma~\ref{lem:support-uniform-lower-cloud} shows that the noncandidate exposed
cloud contributes \(o_{\Pbb,A}(1)\).

Choose \(A_0\) so that \(C_{C,R}Me^{-A_0}\le1/4\), and then take
\(n\ge4M\).  On the lower cloud, the elementary inequality
\[
 |\log(1+v)-v|\le2v^2,
 \qquad v\ge-M/n,\quad |v|\le1/2,
\]
and Lemma~\ref{lem:support-uniform-lower-cloud} give
\[
 \sum_{i:x_{n,i}\le-A}\{\log(1+v_i)-v_i\}
 =O_{\Pbb,C,R,J}(M^2e^{-A})+o_{\Pbb,A}(1).
\]
Since the exposed cloud has \(O_{\Pbb,A}(1)\) observations,
\[
 \sum_{i:x_{n,i}\le-A}v_i
 =\sum_{i:x_{n,i}\le-A}F_{n,i}-m+o_{\Pbb,A}(1).
\]
Equation \eqref{eq:support-lower-cloud-linear} therefore supplies the penalty
\(-m\) in \(d\ge3\), and the penalty
\(-\sum_j\int\{1-B_{\lambda_n}(r)\}\,d\nu_j\) in \(d=2\).
Lemma~\ref{lem:support-uniform-lower-cloud} already includes the anchor error.
Combining the candidate observations, the noncandidate exposed cloud, and the
lower cloud proves \eqref{eq:support-finite-patch-error}.
\end{proof}

An elementary coercivity estimate bounds the edge mass without assuming a
fixed number of fitted atoms; it depends only on the candidate data points,
not on the mixing complexity within their patches.
\begin{proposition}
\label{prop:support-patch-coercivity}
Restrict \eqref{eq:support-anchored-likelihood} to \(J\) edge patches, and let
\(\mathcal I_n\) be their associated observations. Suppose throughout the
patches that
\[
 \frac1n\sum_{i\notin\mathcal I_n}q_{n,\tau,a}(i)\le b<1,
 \qquad
 \max_{i\in\mathcal I_n}\frac1nq_{n,\tau,a}(i)\le C_0.
\]
If \(m=nW\) is the total rescaled edge mass, then, for all large \(n\),
\[
    \mathcal L_{n,\tau}(\nu)
    \le-c_bm+J\log(1+C_0m),
    \qquad c_b=(1-b)/3.
\]
Consequently, every maximizer satisfies
\(m\le4C_0J^2/c_b^2\).
\end{proposition}

\begin{proof}
Fubini's theorem gives
\(\sum_{i\notin\mathcal I_n}F_{n,i}(\nu)\le bm\). Jensen's inequality then
implies
\[
 \sum_{i\notin\mathcal I_n}
 \log\{1-m/n+F_{n,i}(\nu)\}
 \le(n-J)\log\left\{1-\frac mn+\frac{bm}{n-J}\right\}
 \le-c_bm.
\]
At each retained observation the likelihood ratio is at most \(1+C_0m\),
which proves the displayed bound. A maximizer has nonnegative gain, and
\(\log(1+x)\le2\sqrt x\) then yields
\(m\le4C_0J^2/c_b^2\).
\end{proof}

The two-dimensional contact argument requires differentiated, rather than only
uniform, control of the inward-edge background. The next lemma upgrades the
fixed-cutoff decomposition to the \(C^2\) form needed below.
\begin{lemma}
\label{lem:support-d2-c2-absorption}
Assume \(d=2\) and \(\lambda_n\to\lambda\in(0,\infty)\), and put
\(\ell_n^{\rm edge}=4\log\rho_n\). For \(v\in S^1\), define
\[
 a_{n,v}(r,z)=(\rho_n+r)\mathsf R_{z/\rho_n}v
\]
and, for fixed \(R<\infty\), let
\begin{align*}
 \mathcal B_{n,v}(r,z)
 &=\frac1n e^{-\kappa_n|a_{n,v}(r,z)|^2}
   \sum_{i:\,\xi_i<-\ell_n^{\rm edge}}
   e^{a_{n,v}(r,z)\cdot Z_i-|a_{n,v}(r,z)|^2/2},\\
 \mathcal A_{n,v,A}(r,z)
 &=\frac1n e^{-\kappa_n|a_{n,v}(r,z)|^2}
   \sum_{i:\,-\ell_n^{\rm edge}\le\xi_i\le-A}
   e^{a_{n,v}(r,z)\cdot Z_i-|a_{n,v}(r,z)|^2/2}.
\end{align*}
Then
\begin{equation}
\label{eq:support-d2-diffuse-c2}
 \sup_{v\in S^1}
 \|\mathcal B_{n,v}-e^{-\lambda}\Phi(-r)\|_{C^2([-R,R]^2)}
 \xrightarrow{\Pbb}0,
\end{equation}
and, for every \(\delta>0\),
\begin{equation}
\label{eq:support-d2-strip-c2}
 \lim_{A\to\infty}\limsup_{n\to\infty}
 \Pbb\left\{
   \sup_{v\in S^1}\|\mathcal A_{n,v,A}\|_{C^2([-R,R]^2)}
       >\delta
 \right\}=0.
\end{equation}
\end{lemma}

\begin{proof}
Write \(a=a_{n,v}(r,z)\). The radial identity gives
\[
 \frac1n e^{-\kappa_n|a|^2}e^{a\cdot Z_i-|a|^2/2}
 =e^{\xi_i-\kappa_n|a|^2-|Z_i-a|^2/2}.
\]
Cover \(S^1\) by \(O(\rho_n)\) arcs of angular length \(O(1/\rho_n)\), and on
an arc centered at \(v_\alpha\) write
\(v=\mathsf R_{y/\rho_n}v_\alpha\), with \(y\) in a fixed interval. On the
resulting fixed boxes in \((y,r,z)\), differentiation inserts a polynomial of
bounded degree in \(Z_i-\gamma_na\), with coefficients bounded uniformly in
\(n\). Consequently, for every multi-index \(|\beta|\le3\), a summand of the
corresponding derivative of \(\mathcal B_{n,v}\) is bounded by
\(C_Re^{-\ell_n^{\rm edge}}\). Exponential tilting, as in the proof of
Proposition~\ref{prop:two-d-compact-decomposition}, also gives
\[
 \sup_{\alpha,y,r,z}\sum_i
 \E\left|\partial^\beta
   \left[\frac1n e^{-\kappa_n|a|^2}
   e^{a\cdot Z_i-|a|^2/2}
   \mathbf1_{\{\xi_i<-\ell_n^{\rm edge}\}}\right]\right|^2
 \le \frac{C_R}{\rho_n}.
\]
Here \(\partial^\beta\) denotes derivatives in the scaled local variables
\((y,r,z)\), and the supremum ranges over their fixed coordinate boxes.
The corresponding sum of \(p\)th moments is at most
\(C_Re^{-(p-2)\ell_n^{\rm edge}}/\rho_n\). For any fixed \(p>3\),
Rosenthal's inequality therefore bounds the expected \(p\)th power of the
\(W^{3,p}\) norm on one box by
\(C_R\rho_n^{-p/2}+o(\rho_n^{-p/2})\). Summing over the
\(O(\rho_n)\) boxes still gives \(o(1)\). The Sobolev--Morrey embedding
\(W^{3,p}\hookrightarrow C^2\) on each fixed box, followed by Markov's
inequality, proves uniform \(C^2\) concentration over all \(v\in S^1\).

The mean is rotation invariant and exponential tilting gives
\[
 \E\mathcal B_{n,v}(r,z)
 =e^{-\kappa_n(\rho_n+r)^2}
   \Pbb\{|G+(\rho_n+r)e_1|^2
              \le\rho_n^2-2\ell_n^{\rm edge}\},
 \qquad G\sim\mathcal N(0,I_2).
\]
The boundary expansion in the proof of
Proposition~\ref{prop:two-d-compact-decomposition} is uniform after two
derivatives in \(r\); the mean has no \(z\)-dependence. Since
\(\ell_n^{\rm edge}/\rho_n\to0\) and
\(\kappa_n\rho_n^2\to\lambda\), the last display converges in \(C^2\) to
\(e^{-\lambda}\Phi(-r)\). This proves
\eqref{eq:support-d2-diffuse-c2}.

For the strip, let \(Y_i(v)=\rho_n\operatorname{ang}(U_i,v)\). Uniformly on
the same compact set and for \(|\beta|\le2\), the radial identity and the
elementary bound of a polynomial times a Gaussian give
\[
 |\partial_{r,z}^\beta \mathcal A_{n,v,A}(r,z)|
 \le C_R\sum_{i:\,-\ell_n^{\rm edge}\le\xi_i\le-A}
 e^{\xi_i}\{1+|Y_i(v)-z|^4\}e^{-c|Y_i(v)-z|^2}.
\]
Taking the supremum over \(z\in[-R,R]\) only changes the constants and the
Gaussian exponent. The slab-and-arc occupancy proof of
Lemma~\ref{lem:d2-low-atom-absorption} applies with the summable weight
\((1+q^4)e^{-cq^2}\) in place of \(e^{-q^2/2}\). It bounds the last display
by \(C_Re^{-A}+o_{\Pbb}(1)\), simultaneously for \(|\beta|\le2\), and proves
\eqref{eq:support-d2-strip-c2}.
\end{proof}

To convert each local maximizer into an exact contact count, we need
two-derivative convergence of the fitted KKT certificate; convergence of
likelihood values alone would not exclude nearby split atoms.
\begin{proposition}
\label{prop:support-one-contact-per-patch}
Suppose an NPMLE is localized to finitely many compact patches,
its rescaled edge mass is bounded, and every associated radial excess is at
least \(\eta>0\) from the relevant threshold. After the radial cutoffs are
removed, an inactive patch contains no support point and an active patch
contains exactly one, apart from an event of probability tending to zero.
Here active means \(x_{n,j}>s_n\) for \(d\ge3\) and
\(x_{n,j}>x_{\lambda_n}\) for \(d=2\).
\end{proposition}

\begin{proof}
It is enough to argue along subsequences. The radial cutoffs and threshold
buffer allow a further subsequence on which the number of candidate patches is
constant and all marks \(x_{n,j}\) converge, say to \(x_j\). Write
\(a_{n,j}(u)=a_{n,j}^{(\ge3)}(u)\) in \(d\ge3\), and
\(a_{n,j}(u)=a_{n,j}^{(2)}(r,z)\) for \(u=(r,z)\) in \(d=2\).
The candidate index set is fixed throughout this subsequence; only the lower
observational cutoff \(A\) will later vary.  For a localized measure
\(\nu=(\nu_1,\ldots,\nu_J)\), let
\[
 M_{n,i}(\nu)=1+\frac1n\sum_j
   \int\{q_{n,\tau,a_{n,j}(u)}(i)-1\}\,\nu_j(du).
\]
Moving mass from the central atom into patch \(j\) has derivative
\begin{equation}
\label{eq:support-finite-patch-certificate}
 \Gamma_{n,j,\nu}(u)
 =\frac1n\sum_{i=1}^n
   \frac{q_{n,\tau,a_{n,j}(u)}(i)-1}{M_{n,i}(\nu)}.
\end{equation}
It is nonpositive throughout the patch and vanishes at every support point.

The limiting certificate in an active \(d\ge3\) patch is
\[
    \Gamma_j^*(h)=e^{-|h|^2/2}-1;
\]
it has a unique zero at zero and Hessian \(-I_d\). In an active
two-dimensional patch it is
\[
 \Gamma_j^*(r,z)
 =(1-B_\lambda(r))
 \left\{\frac{R_{x_j}(r,z)}{\sup_{u,y}R_{x_j}(u,y)}-1\right\},
\]
whose unique zero is \((-a_\lambda,0)\) and whose Hessian there is
\(-\{1-e^{-\lambda}\Phi(a_\lambda)\}I_2\). A threshold buffer \(\eta\)
also gives a fixed negative certificate gap in every inactive patch.

We next verify that \eqref{eq:support-finite-patch-certificate} converges to
these limits in \(C^2\) on a smaller patch, uniformly over measures in a fixed
compact mass ball.  Fix the observational cutoff \(A\).  At the candidate
observations, completing the square and differentiating three times gives
uniform \(C^3\) convergence of the associated response.  Every response from
a different exposed observation carries the factor
\(e^{-c(\ell_n^{\rm sep})^2}\), and the exposed cloud has tight cardinality.

For the lower observations in fixed \(d\ge3\),
Lemma~\ref{lem:support-dge3-palm-c2} gives, simultaneously in all candidate
patches, \(C^2\) convergence of the punctured lower field to zero.  The lemma
is centered at \(Z_j/\gamma_n\), whereas the score-displacement patch is
centered at \(Z_j/\gamma_n-\tau\). Taking the Palm estimate on a slightly
larger compact set and using \(|\tau|=o_{\Pbb}(1)\) transfers the bound to the
score-displacement patch. The accompanying multiplier
\(e^{-\gamma_na\cdot\tau}\) is \(1+o_{\Pbb}(1)\) in \(C^2\), because
\(\rho_n|\tau|=o_{\Pbb}(1)\).  Thus this is a deterministic translation of an
already uniform \(C^2\) estimate and requires no conditioning on a
sample-centered exceedance set.

In \(d=2\), Lemma~\ref{lem:support-d2-c2-absorption} gives \(C^2\)
convergence of the diffuse background to \(e^{-\lambda}\Phi(-r)\) and removes
the strip between the receding and fixed radial cutoffs in the \(C^2\)
topology.
The estimate is uniform over the patch direction, so it remains valid for the
random directions \(U_j\). The KKT field differs from the field in
Lemma~\ref{lem:support-d2-c2-absorption}
by the anchor multiplier
\[
 e^{-a\cdot\Delta_n^{\rm anc}-a\cdot\bar Z_n}
 =e^{-\gamma_na\cdot\tau}.
\]
Lemma~\ref{lem:support-fixed-core-collapse} gives
\[
 \rho_n^2|\tau|^2
 =O_{\Pbb}\left\{\frac{(\log n)^{C_d+1}}{n\kappa_n}\right\}
 =o_{\Pbb}(1),
\]
and differentiation of the multiplier in the local coordinates only
introduces further factors \(O_{\Pbb}(|\tau|)\).
It is therefore \(1+o_{\Pbb}(1)\) in \(C^2\) on every fixed patch.

It remains to pass from the unfitted fields to the fitted certificate.  Take
\(A\ge A_0(C,R,M)\) from
Lemma~\ref{lem:support-uniform-lower-cloud}.  For every lower observation,
\[
 |M_{n,i}(\nu)-1|
 \le F_{n,i}(\nu)+M/n
 \le C_{C,R}Me^{-A}+o(1)\le\frac12,
\]
uniformly over measures of mass at most \(M\).  Hence
\(|M_{n,i}^{-1}-1|\le2|M_{n,i}-1|\).  Applying
Lemma~\ref{lem:support-uniform-lower-cloud} after each of
the at most two local differentiations gives
\begin{equation}
\label{eq:support-weighted-lower-c2}
 \sup_{\nu:\,\nu(\mathbb R^d)\le M}
 \left\|
  \frac1n\sum_{i:x_{n,i}\le-A}
 \{q_{n,\tau,a_{n,j}(\cdot)}(i)-1\}
 \{M_{n,i}(\nu)^{-1}-1\}
 \right\|_{C^2}
 =O_{\Pbb,C,R,J}(Me^{-A})+o_{\Pbb,A}(1),
\end{equation}
with a tight \(O_{\Pbb}(1)\) coefficient uniform for \(A\ge A_0\).
To see this directly, note that \(M_{n,i}(\nu)\) is fixed when the prospective
contact location \(u\) is varied.  Thus all derivatives fall on the numerator,
and \eqref{eq:support-lower-cloud-derivative-envelope}, together with
\(\max_i|M_{n,i}^{-1}-1|\le CMe^{-A}+o(1)\), gives the displayed bound.  The
constant \(-1\) at derivative order zero is controlled by the same maximum.
At a candidate observation, the denominator converges to its local-program
denominator.  At a noncandidate exposed observation, every candidate response
is \(o_{\Pbb,A}(1)\), uniformly in \(C^2\); the constant \(-1\) in the
numerator contributes only \(O_{\Pbb,A}(1/n)\) after summing over the tight
exposed cloud.  Thus the complete fitted certificate
converges sequentially, first as \(n\to\infty\) and then as \(A\to\infty\),
to the displayed limiting certificates.

We now make the order of the cutoff and argmax limits explicit. Proposition
\ref{prop:support-patch-coercivity} confines the total rescaled mass to a
deterministic compact interval on events of arbitrarily high probability.
Together with the fixed compact candidate-patch union, this gives a weakly
compact space \(\mathcal M_{M,R}\) of candidate measures. Let
\(\mathcal J^\star_{\boldsymbol x}\) be the sum of the limiting local programs
for a vector \(\boldsymbol x=(x_1,\ldots,x_J)\) of candidate marks, and let
\(\nu^\star_{\boldsymbol x}\) be its optimizer. The scalar optimizations above
show that this optimizer is unique and depends continuously on
\(\boldsymbol x\) as long as the marks remain outside the threshold buffer.

Choose the patch radius \(R\) so that every such optimizer is in the interior.
For a fixed small \(\delta>0\), let
\(\mathcal U_{\boldsymbol x}\) be its radius-\(\delta\) weak neighborhood in
\(\mathcal M_{M,R}\).  The mark vectors with \(|x_j|\le C\) and distance at
least \(\eta\) from the threshold form a compact set.  Joint continuity of
the local programs and uniqueness of their optimizers therefore give the
uniform separation
\begin{equation}
\label{eq:support-argmax-gap}
 c_{\delta,\eta,C,M,R,J}
 :=\inf_{\boldsymbol x}
 \left[
  \mathcal J^\star_{\boldsymbol x}(\nu^\star_{\boldsymbol x})
  -\sup_{\nu\in\mathcal M_{M,R}\setminus\mathcal U_{\boldsymbol x}}
      \mathcal J^\star_{\boldsymbol x}(\nu)
 \right]>0.
\end{equation}
Decrease \(\delta\), if necessary, so that every active patch has mass bounded
away from zero throughout \(\mathcal U_{\boldsymbol x}\).  All the objects in
\eqref{eq:support-argmax-gap} depend only on the threshold-buffered candidate
patches, not on the observational cutoff.  Choose the deterministic cutoff
\(A\ge A_0(C,R,M)\) so large that, by
\eqref{eq:support-finite-patch-iterated}, the finite-patch approximation error
exceeds one quarter of the gap in \eqref{eq:support-argmax-gap} with
arbitrarily small limiting probability.  Then take \(n\) large.  Uniform
convergence of the local kernels shows on the complementary event that every
finite-sample maximizer belongs to \(\mathcal U_{\boldsymbol x}\): otherwise
comparison with the pushforward of \(\nu^\star_{\boldsymbol x}\) would lose at
least half that gap.  In particular, every active finite-sample patch carries
positive mass and hence contains a KKT contact.

Increase this same deterministic \(A\), if necessary, so that
\eqref{eq:support-weighted-lower-c2} exceeds one quarter of either the
inactive-patch certificate gap or the absolute value of the largest active
limiting Hessian eigenvalue with arbitrarily small limiting probability.
These constants are uniform over the threshold-buffered compact mark set,
while the candidate space does not change.  An inactive patch then has no
contact.  In an active patch,
\(C^0\) convergence confines all contacts to a small convex neighborhood of
the unique limiting zero, and \(C^2\) convergence makes the finite-sample
certificate strictly concave there.  Since the certificate is nonpositive,
strict concavity rules out two distinct zeros.  Sending the two exceptional
probabilities to zero proves the claim along every subsequence.
\end{proof}

\subsubsection{Threshold pruning and the support law}

Only observations within a fixed radial distance of the activation threshold
can support atoms. The needed deleted-field estimate is a further consequence
of the Bennett--Morrey argument in
Section~\ref{sec:sharp-high-dimensional-estimates}.
\begin{lemma}
\label{lem:support-dge3-deleted-field}
Assume \(d\ge3\) and \(s_n=O(1)\). For fixed \(a>0\), define
\[
 S_t^{(a)}=\frac1n e^{-\kappa_n|t|^2}
 \sum_{i=1}^ne^{t\cdot Z_i-|t|^2/2}
 \mathbf1_{\{|Z_i|^2/2\le\log n+\lambda_n-a\}}.
\]
If \(\zeta\in(0,1)\) and \(e^a\zeta>1\), then
\[
    \Pbb\left\{\sup_{|t|\ge\rho_n/2}S_t^{(a)}\ge\zeta\right\}\to0.
\]
\end{lemma}

\begin{proof}
Put \(R_a^2/2=\log n+\lambda_n-a\) and
\(h=R_a-|t|\). On each fixed physical annulus \(|h|\le H\),
Lemma~\ref{lem:physical-offset-fixed-t} gives
\[
 \Pbb\{S_t^{(a)}\ge\zeta\}
 \le\left(\rho_n^{-1}e^{-\lambda_n}\right)^{
     e^{a+h^2/2}\zeta-o(1)}.
\]
Lemma~\ref{lem:scaled-derivative-morrey} reduces the annulus to a net of
cardinality \(\rho_n^{d-1+o(1)}\), with \(o(1)\) oscillation. Since
\(\rho_n^{-1}e^{-\lambda_n}=\rho_n^{-(d-1)+o(1)}\) and
\(e^a\zeta>1\), the union bound tends to zero.

For \(h\ge H\), Lemma~\ref{lem:growing-offset-fixed-t} gives
\[
 \Pbb\{S_t^{(a)}\ge\zeta/2\}
 \le\exp\{-c_de^{a+h^2/4}\log\rho_n\}.
\]
A polynomial-size Euclidean net and the same derivative control make the
union bound vanish after \(H\) is large. When \(|t|\ge R_a+H\), every
truncated summand decreases along its radial ray, so the outward region reduces
to the controlled boundary. These ranges cover \(|t|\ge\rho_n/2\).
\end{proof}

The two-dimensional analogue of the deleted-field estimate has an explicit
profile gap.  It is useful to record both the global residual bound and the
stronger bound on the candidate patches, where the associated extreme
observation has itself been deleted.
\begin{lemma}
\label{lem:support-d2-residual-gap}
Assume \(d=2\) and \(\lambda_n\to\lambda\in(0,\infty)\).  Fix \(\eta>0\),
put
\[
 \mathcal I_n(\eta)=\{i:\xi_i>x_{\lambda_n}-\eta\},
 \qquad
 \mathcal R_{n,\eta}(a)
 =\frac1n\sum_{i\notin\mathcal I_n(\eta)}q_{n,\tau,a}(i),
\]
and let \(\mathcal P_n(\eta,R)\) be the union of the radius-\(R\) local
patches associated with \(\mathcal I_n(\eta)\).  Then
\begin{align}
\label{eq:support-d2-global-residual-gap}
 \sup_{|a|\ge\rho_n/2}\mathcal R_{n,\eta}(a)
 &\le1-\frac12(1-e^{-\eta})(1-e^{-\lambda})+o_{\Pbb}(1),\\
\label{eq:support-d2-candidate-residual-gap}
 \sup_{a\in\mathcal P_n(\eta,R)}\mathcal R_{n,\eta}(a)
 &\le e^{-\lambda}+o_{\Pbb,R}(1)
\end{align}
for each fixed \(R\).
\end{lemma}

\begin{proof}
Fix \(C,K,A<\infty\) and work first on
\(\max_i\xi_i\le C\) and \(|\,|a|-\rho_n|\le K\).  Write
\(a=(\rho_n+r)v\), \(|r|\le K\).  Lemma
\ref{lem:d2-fixed-cutoff-decomposition}, its quantitative absorption bound,
and the anchor estimate give, uniformly in \(r\) and \(v\),
\begin{equation}
\label{eq:support-d2-residual-decomposition}
 \mathcal R_{n,\eta}((\rho_n+r)v)
 =B_{\lambda_n}(r)
  +\sum_{\substack{-A<\xi_i\le x_{\lambda_n}-\eta}}
    C_{\lambda_n,\xi_i}(r,Y_i(v))
  +O_{\Pbb,K}(e^{-A})+o_{\Pbb,A,K}(1).
\end{equation}
For fixed \(A,C\), the exposed points have tight cardinality and, with
probability tending to one, their pairwise scaled angular distances exceed a
deterministic \(\ell_n\to\infty\).  If there are at most \(J\) exposed points,
then for any \(v\) all but at most one have \(|Y_i(v)|\ge\ell_n/2\); uniformly
for \(|r|\le K\), their total contribution is at most
\(J e^{C+K^2/2}e^{-\ell_n^2/8}\).  Tightness of \(J\) therefore makes this
remainder \(o_{\Pbb,A,C,K}(1)\), uniformly in \(v\).  Since every retained
residual mark is at most
\(x_{\lambda_n}-\eta\),
\begin{align*}
 B_{\lambda_n}(r)
 +C_{\lambda_n,x_{\lambda_n}-\eta}(r,z)
 &=e^{-\eta}
   \{B_{\lambda_n}(r)+C_{\lambda_n,x_{\lambda_n}}(r,z)\}
   +(1-e^{-\eta})B_{\lambda_n}(r)\\
 &\le e^{-\eta}+(1-e^{-\eta})e^{-\lambda_n}\\
 &=1-(1-e^{-\eta})(1-e^{-\lambda_n}).
\end{align*}
Choose \(A\) so that the absorption error is less than one quarter of the
limiting gap and then take \(n\to\infty\).  This proves the first bound on
every fixed radial window, on \(\{\max_i\xi_i\le C\}\).  As \(|r|\to\infty\),
the spike term is bounded by \(e^{C-r^2/2+o(1)}\), while
\(B_{\lambda_n}(r)\) tends to \(e^{-\lambda}\) inward and to zero outward.
The diffuse- and tail-field estimates in the proof of
Proposition~\ref{prop:d-two-radial-compactness} make these bounds uniform as
\(K\to\infty\).  Finally,
\[
 \limsup_{n\to\infty}\Pbb\{\max_i\xi_i>C\}
 =1-\exp\{-e^{-C}\}\longrightarrow0
 \qquad(C\to\infty),
\]
which removes the upper cutoff and proves
\eqref{eq:support-d2-global-residual-gap}.

If \(a\in\mathcal P_n(\eta,R)\), its associated observation belongs to
\(\mathcal I_n(\eta)\) and is absent from the residual sum.  Every exposed
residual direction is then at scaled angular distance at least
\(\ell_n-R\), so the entire exposed sum in
\eqref{eq:support-d2-residual-decomposition} is \(o_{\Pbb,A,C,R}(1)\).
The remaining profile is at most
\(\sup_rB_{\lambda_n}(r)=e^{-\lambda_n}\).  Taking first \(n\to\infty\),
then \(A,C\to\infty\), proves
\eqref{eq:support-d2-candidate-residual-gap}.
\end{proof}

Combining the edge and bulk gaps prunes subthreshold points before the
mass-coercivity bound is used. Thus the number of candidate patches is tight
independently of the lower radial cutoff.
\begin{proposition}
\label{prop:support-threshold-pruning}
Fix \(\eta>0\), and set
\[
    \vartheta_n=
    \begin{cases}
       s_n,&d\ge3,\\
       x_{\lambda_n},&d=2.
    \end{cases}
\]
In the corresponding critical window, every noncentral support point of an
NPMLE lies, with probability tending to one after the radial
cutoffs are removed, in a patch of some deterministic radius
\(R=R(\eta)<\infty\) associated with an observation whose radial excess
satisfies \(x_{n,i}>\vartheta_n-\eta\). The number of these patches is
\(O_{\Pbb}(1)\). More precisely, put
\[
 \mathcal I_n(\eta)=\{i:x_{n,i}>\vartheta_n-\eta\},
\]
and let \(\mathcal P_n(\eta)\) denote the union of their radius-\(R\) patches.
There are a deterministic \(b<1\) and
random variables \(C_{0,n}=O_{\Pbb}(1)\) such that
\[
 \sup_{a\in\mathcal P_n(\eta)}
 \frac1n\sum_{i:x_{n,i}\le\vartheta_n-\eta}
 q_{n,\tau,a}(i)\le b,
 \qquad
 \sup_{a\in\mathcal P_n(\eta)}
 \max_{i:x_{n,i}>\vartheta_n-\eta}
 \frac1nq_{n,\tau,a}(i)\le C_{0,n}.
\]
\end{proposition}

\begin{proof}
Call the sum over \(i\notin\mathcal I_n(\eta)\) the residual field. For
\(d\ge3\), every residual
observation satisfies
\[
    |Z_i-\bar Z_n|^2/2\le b_{n,d}+s_n-\eta
       =\log n+\lambda_n-\eta.
\]
Uniformly in \(i\),
\[
 \frac{|Z_i|^2}{2}
 \le\frac{|Z_i-\bar Z_n|^2}{2}
    +\max_i|Z_i-\bar Z_n|\,|\bar Z_n|
    +\frac{|\bar Z_n|^2}{2}
 =\frac{|Z_i-\bar Z_n|^2}{2}+o_{\Pbb}(1).
\]
Thus, with probability tending to one, every residual observation lies below
\(\log n+\lambda_n-\eta/2\). Put
\[
    a_\eta=\eta/2,
    \qquad
    \zeta_\eta=\frac{1+e^{-a_\eta}}2.
\]
Since \(e^{a_\eta}\zeta_\eta>1\), Lemma
\ref{lem:support-dge3-deleted-field} applied with \(a=a_\eta\) gives a gap of
\((1-\zeta_\eta)/2\) below one throughout
\(|a|\ge\rho_n/2\), with probability tending to one. To pass from the raw
field in Lemma~\ref{lem:support-dge3-deleted-field} to the residual field, note that
 \[
 \frac1n\sum_{\rm residual}q_{n,\tau,a}(i)
 \le e^{-a\cdot\Delta_n^{\rm anc}-a\cdot\bar Z_n}S_a^{(a_\eta)}.
 \]
An NPMLE may be supported in the convex hull of the rescaled observations, so
\(|a|=O_{\Pbb}(\rho_n)\); on this range
\eqref{eq:support-high-d-anchor-error} and
\(\rho_n|\bar Z_n|=o_{\Pbb}(1)\) make the prefactor \(1+o_{\Pbb}(1)\).
Thus the residual field is at most
\((1+\zeta_\eta)/2<1\) throughout the \(d\ge3\) edge region.

In \(d=2\), Lemma~\ref{lem:support-d2-residual-gap} gives the explicit global
edge gap
\[
 g_{2,\eta}=\frac12(1-e^{-\eta})(1-e^{-\lambda})>0.
\]
On the candidate patches themselves it gives the stronger residual bound
\(e^{-\lambda}+o_{\Pbb}(1)\). Thus, in the first displayed inequality of the
proposition, one may take any fixed \(b\in(e^{-\lambda},1)\) for all
sufficiently large \(n\).

The bulk estimates in Proposition~\ref{prop:bulk-mgf-exclusion} for
\(d\ge3\), and Lemma~\ref{lem:d2-bulk-exclusion} for \(d=2\), retain a gap
of order \(\kappa_n\) on the region between the fixed core and
\(\rho_n/2\). The candidate set has tight cardinality by radial Poisson
convergence, and its maximum radial excess is tight. On an event of arbitrarily
high probability, both its cardinality and its normalized peak heights are
bounded by deterministic constants. The exact square completion
\eqref{eq:support-lower-cloud-square-completion} then bounds the total
candidate contribution outside radius \(R\) of the patch centers by
\(Ce^{-cR^2}\). Choose \(R\) so this is smaller than one quarter of the
relevant residual edge gap: \((1-\zeta_\eta)/2\) for \(d\ge3\), and
\(g_{2,\eta}\) for \(d=2\). The full central field consequently has a fixed
gap outside \(\mathcal P_n(\eta)\) throughout the edge region. Equation
\eqref{eq:support-high-d-anchor-error} transfers both gaps to the central
anchor. Finally, Lemma
\ref{lem:support-no-percolation} preserves the edge gap because
\(W_n=o_{\Pbb}(1)\) and the bulk gap because
\(W_n=o_{\Pbb}(\kappa_n)\). KKT equality therefore excludes every
support point outside the candidate patches.

On the fixed candidate patches,
Lemma~\ref{lem:support-dge3-deleted-field} for \(d\ge3\) and
Lemma~\ref{lem:support-d2-residual-gap} for \(d=2\) give the first displayed
inequality with deterministic \(b<1\). The tight upper radial cutoff,
square completion, and angular separation give the second with
\(C_{0,n}=O_{\Pbb}(1)\).
\end{proof}

We can now count the patches. The following buffered sandwich is the
fixed-dimensional input to the common edge-count reduction.
\begin{proposition}
\label{prop:fixed-dimensional-edge-sandwich}
Assume \(\lambda_n\to\lambda\in(0,\infty)\) when \(d=2\), or
\(s_n^\star\to s\in\mathbb R\) when \(d\ge3\) is fixed. With
\(x_{n,i}\) defined in \eqref{eq:fixed-support-radial-excess} and
\(\vartheta_n\) as in Proposition~\ref{prop:support-threshold-pruning}, for
every fixed \(\eta>0\),
\begin{equation}
\label{eq:support-count-sandwich}
 \#\{i:x_{n,i}>\vartheta_n+\eta\}
 \le\widehat K_n-1
 \le\#\{i:x_{n,i}>\vartheta_n-\eta\},
\end{equation}
apart from an event whose limiting probability is \(O(\eta)\).
Under either critical-window assumption, the NPMLE is unique with probability
tending to one.
\end{proposition}

\begin{proof}
First assume \(\lambda_n\to\lambda\in(0,\infty)\) when \(d=2\), or
\(s_n^\star\to s\in\mathbb R\) when \(d\ge3\). Fix \(\eta>0\). Proposition
\ref{prop:support-threshold-pruning} confines every noncentral support point
to the \(O_{\Pbb}(1)\) patches associated with
\(x_{n,i}>\vartheta_n-\eta\), and supplies the two hypotheses of Proposition
\ref{prop:support-patch-coercivity}. Hence the rescaled edge mass is
\(O_{\Pbb}(1)\). After choosing an upper radial cutoff and a deterministic
bound on the number of patches, Proposition
\ref{prop:support-finite-patch-expansion} and Proposition
\ref{prop:support-one-contact-per-patch} show that every active patch contains
exactly one fitted atom and every inactive candidate patch contains none.
These are precisely the inequalities in \eqref{eq:support-count-sandwich}. Removing
the radial cutoffs costs \(o(1)\); the only remaining exceptional event is a
radial point in the threshold buffer, whose probability is \(O(\eta)\).

On the event where Propositions~\ref{prop:support-threshold-pruning} and
\ref{prop:support-one-contact-per-patch} give the stated patch localization,
the NPMLE is also unique. If \(G\) and \(H\) are two NPMLEs, then
\((G+H)/2\) is also an NPMLE.  Lemma~\ref{lem:support-fixed-core-collapse},
applied to all three measures, forces their central atoms to coincide:
otherwise the central part of \((G+H)/2\) would have positive covariance.
Propositions~\ref{prop:support-threshold-pruning} and
\ref{prop:support-one-contact-per-patch}, applied to \((G+H)/2\), likewise
force the edge contact in every active patch to coincide.  Thus all NPMLEs
are supported on the same finite contact set.

The likelihood vectors at these contacts are linearly independent.  Indeed,
on the rows indexed by the active observations, the edge-contact block has
diagonal entries of order \(n\), bounded away from zero after division by
\(n\), and off-diagonal entries \(o(n)\).  The finite-patch estimates also
give an \(O_{\Pbb}(1)\) empirical average for each of the finitely many edge
columns.  Hence there is a noncandidate row on which every edge entry is
\(o(n)\); adjoining this row and the central column produces, after scaling
the edge columns by \(n^{-1}\), an asymptotically invertible triangular
minor.  Strict concavity of the log likelihood makes the fitted likelihood
vector common to all maximizers, and linear independence then makes its
representing weights unique.  The probability of the exceptional event is
\(O(\eta)+o(1)\); letting \(\eta\downarrow0\) proves the uniqueness claim.
\end{proof}
